\section{Conclusion}
\label{sec:limits}

This paper has asked whether a symbolic-regression search can recover a structural equation when the regressors are endogenous, and has answered by moving the instrument from the generation of candidate formulas to their scoring. Candidates may come from any cheap generator, a squared-error search, a distilled nonparametric IV curve, a sparse library path or a list of textbook shapes; each is refitted by GMM on sieve instruments, the formula is selected by a GMM model-selection criterion with a per-node structure penalty, and an Anderson--Rubin confidence set over formulas reports whether the data identify the result. The approach treats symbolic regression and instrumental variables as complements: the search supplies the functional forms that econometrics would otherwise have to assume, and the moment conditions supply the identification that the search lacks.

The theoretical results characterize when this works. The selection rule is consistent for any generator whose candidates contain the true form, the constants of the selected formula converge at the parametric rate, and split-sample GMM inference on it is valid (\res{Theorem}{1}). The finite-sample bound of \res{Theorem}{2} isolates the two channels through which selection fails, a wrong formula that the moments barely separate from the truth and a longer formula that fits the moment noise, and shows that the structure penalty is what controls the second channel uniformly in the formula size. \res{Proposition}{5} establishes that the formula confidence set covers the true form without any rank condition, which is what allows the method to report weak identification rather than a confident formula, and \res{Proposition}{4} shows that a $J$-point instrument cannot exclude formulas with $J$ or more constants, so that discrete instruments bound the complexity the moments can discipline.

The Monte Carlo evidence is consistent with these results. Under confounding and a heteroskedastic first stage the method recovers the log, power and two-regressor laws in 92--100\% of runs, and in 60--92\% under simultaneity and a discrete instrument; in the strong-endogeneity setting its recovery of the log and power laws reaches 98--100\% from $n = 5000$, whereas the squared-error search never recovers the power law under strong endogeneity and, where it finds the log form, carries a slope bias equal to $\rho/(1 + \pi^2)$ to two decimals. The contribution of the instrument is concentrated in the scoring step: candidates from a plain search, from a distilled sieve-2SLS curve or from a library of economic templates, selected by the moment criterion, perform as well as or better than a search run on the moment criterion itself. Under a weak instrument recovery falls to 0--2\% and the diagnostic flags every run, at a 4\% false-alarm rate under strong instruments. On the 1995 Family Expenditure Survey the method reproduces the Working--Leser food share and the quadratic motor share of the Engel-curve literature, and on the quarter-of-birth instruments of \citet{angristkrueger1991} it reports weak identification rather than a formula.

Two limitations of the method are documented in Sections~\ref{sec:mc}--\ref{sec:budget} and not resolved. A threshold law whose kink the instruments barely move is not recovered at $n = 2000$, because a simpler formula fits the moments equally well, and a backward-bending labour supply whose bend lies at the edge of the data is found by a large search in 90\% of runs but selected in 5\%, for the same reason. Both are failures of identification within the formula class rather than of the selection rule, and the diagnostic flags only part of them, since a wrong formula that fits the moments is, by the data's lights, identified. A larger search budget does not help in either case; instruments whose support extends through the kink or the bend, or substantially larger samples, would.

Several questions remain open. Rates of convergence for a growing formula size reduce, by \res{Proposition}{3.3}, to a modulus of continuity on the formula sieve; because the sieve less $f_0$ is not a linear space, the modulus can fall below any quadratic, and bounding it for specific sieves, free-knot splines first, is the natural next step. \res{Theorem}{3(b)} is stated on a compact support; the lognormal designs require a growth envelope in place of the Lipschitz ball, approximants that respect it (clipping with \code{max} and \code{min} does, at the cost of nodes), and a concentration inequality under an unbounded envelope, each of which appears routine and none of which is written out. \res{Theorem}{2} covers the sieve criterion with any fixed positive definite weight; the exact kernel criterion is a degenerate V-statistic at the truth and requires concentration for degenerate U-processes, and an estimated efficient weight enters through \res{Proposition}{2.4} and requires a concentration bound for the weight estimate. For discrete instruments, whether a wrong tree with fewer than $J$ constants is generically identified is a transversality question, stated as \res{Conjecture}{4(c)} in Appendix~\ref{app:prop4}, whose regular-value hypothesis we cannot verify in general. The consistency result is pointwise in the data-generating process; a uniform version over designs whose separation is bounded below would follow from \res{Theorem}{2} with explicit sample sizes. Finally, the additive error excludes random coefficients, and the IV quantile model of \citet{chernozhukovhansen2005} is the natural route to a non-separable extension.

For applied work the results suggest a simple protocol. Candidates should be generated cheaply and include templates for the shapes that economic theory names; the formula should be refitted and selected by the moment criterion with the BIC penalty per constant and the structure term per node; the number of constants should be capped at $J - 1$ when the instrument takes $J$ values; inference should use a split sample; and the identification diagnostic should be reported in every case, since it is the method's answer where the data do not identify a formula. A larger search is optional, and its candidates should be pruned before they are pooled. The code, the run records and a replication guide accompany the paper.
